\documentclass[12pt]{article}

\usepackage{sbc-template}
\usepackage{graphicx,url}
\usepackage[brazil]{babel}
\usepackage[utf8]{inputenc}
\usepackage{listings}
\usepackage{amsmath}
\usepackage{booktabs}

\sloppy

% ============================================================
%  TÍTULO E AUTORES
% ============================================================
\title{Navegação Indoor Inteligente com Fusão de\\
       Beacons BLE e Dead Reckoning Pedestre via Filtro de Madgwick}

\author{
  Murillo Pires Teixeira Melo\inst{1},
  João Pedro Lucena Lisboa Borges\inst{1},
  Pedro Lukas Alves de Sousa Santos\inst{1}
}

\address{
  Instituto de Informática -- INF -- Universidade Federal de Goiás (UFG)\\
  Goiânia -- Goiás -- Brasil\\
  \email{\{murillomelo, borges.joao, pedrolukas\}@discente.ufg.br}
}

% ============================================================
\begin{document}

\maketitle

% ============================================================
%  ABSTRACT (inglês)
% ============================================================
\begin{abstract}
The unavailability of GPS signals in large enclosed environments such as shopping
centers, hospitals, and university campuses creates significant challenges for
pedestrian navigation. Existing solutions often rely on costly infrastructure or
proprietary hardware, limiting large-scale adoption. This paper presents an
open-source indoor navigation system that fuses Bluetooth Low Energy (BLE) beacon
anchoring with Pedestrian Dead Reckoning (PDR) using IMU sensor fusion via the
Madgwick filter. The system adopts an offline-first architecture, synchronizing
a versioned geospatial graph to the mobile device and computing accessible routes
via pgRouting. The tool is fully available under a CC BY-SA 4.0 license, including
a containerized deployment via Docker, and was evaluated in a real indoor scenario
at UFG's Institute of Informatics.
\end{abstract}

% ============================================================
%  RESUMO (português)
% ============================================================
\begin{resumo}
A indisponibilidade do sinal GPS em ambientes fechados de grande escala, como
\textit{shoppings}, hospitais e \textit{campi} universitários, impõe desafios
significativos à navegação pedestre. Soluções existentes frequentemente dependem
de infraestrutura custosa ou hardware proprietário, limitando sua adoção em escala.
Este artigo apresenta um sistema de navegação indoor de código aberto que combina
ancoragem por beacons Bluetooth Low Energy (BLE) com Dead Reckoning Pedestre (PDR)
por meio de fusão de sensores IMU utilizando o Filtro de Madgwick. O sistema adota
arquitetura \textit{offline-first}, sincronizando um grafo geoespacial versionado
para o dispositivo móvel e calculando rotas acessíveis via pgRouting. A ferramenta
está disponível integralmente sob licença CC BY-SA 4.0, com implantação conteinerizada
via Docker, e foi avaliada em cenário indoor real no Instituto de Informática da UFG.
\end{resumo}

% ============================================================
\section{Introdução}
% ============================================================

Ambientes construídos de grande escala — \textit{shoppings}, hospitais, aeroportos
e \textit{campi} universitários — concentram diariamente milhões de usuários que
dependem de orientação espacial eficiente. Estima-se que 80\% do tempo das pessoas
seja passado em ambientes internos \cite{Zafari2019}, onde o GPS apresenta degradação
severa de desempenho em razão da atenuação e do multipercurso dos sinais satelitais
\cite{Liu2007}. Essa lacuna tecnológica afeta de modo desigual usuários com mobilidade
reduzida, que dependem de informações de acessibilidade em tempo real para planejar
seus percursos.

As abordagens existentes para posicionamento indoor podem ser agrupadas em três
categorias: baseadas em infraestrutura dedicada (\textit{Ultra-Wideband}, câmeras),
baseadas em sinais de oportunidade (Wi-Fi, BLE) e baseadas em sensores inerciais
(PDR/IMU). Soluções puramente infraestruturais oferecem alta precisão, porém com
custo de implantação proibitivo. Soluções puramente inerciais, por sua vez, acumulam
erro de deriva que inviabiliza uso prolongado sem mecanismos de correção externos
\cite{Jimenez2009}.

Este trabalho propõe e implementa uma \textbf{ferramenta de navegação indoor} que
endereça essas limitações por meio de uma abordagem híbrida de baixo custo: beacons
BLE comerciais fornecem âncoras de posição esporádicas, enquanto o Filtro de Madgwick
aplicado sobre os sensores IMU do próprio \textit{smartphone} mantém o rastreamento
contínuo entre âncoras. O roteamento é computado sobre um grafo geoespacial persistido
em PostgreSQL com as extensões PostGIS e pgRouting, com suporte nativo a perfis de
acessibilidade. A arquitetura \textit{offline-first} elimina dependência de conectividade
contínua durante a navegação.

As principais contribuições deste trabalho são:
\begin{itemize}
  \item Uma arquitetura de fusão sensorial BLE + PDR de baixo custo, sem hardware
        dedicado além de beacons comerciais;
  \item Um \textit{backend} RESTful com roteamento acessível e sincronização versionada
        por delta, viabilizando operação \textit{offline};
  \item Uma ferramenta completa e reproduzível, distribuída com contêineres Docker
        e licença aberta (CC BY-SA 4.0);
  \item Analytics de fluxo de pessoas em conformidade com a LGPD, destinado a
        gestores de espaços físicos.
\end{itemize}

% ============================================================
\section{Trabalhos Relacionados}
% ============================================================

A literatura de posicionamento indoor tem sido extensivamente revisada nos últimos
anos \cite{Zafari2019, Liu2007}. As abordagens podem ser sistematizadas em função
de três critérios: precisão, custo de infraestrutura e dependência de conectividade.

Sistemas baseados exclusivamente em RSSI de beacons BLE — como iBeacon (Apple) e
Eddystone (Google) — atingem precisão de 1 a 3 metros em condições controladas,
porém sofrem degradação acentuada em ambientes com obstáculos e tráfego humano
intenso \cite{Faragher2015}. Técnicas de \textit{fingerprinting} Wi-Fi melhoram a
robustez em ambientes densos, mas exigem mapeamento prévio extenso e são sensíveis
a mudanças no ambiente.

Soluções baseadas em PDR/IMU oferecem continuidade de rastreamento sem infraestrutura,
mas acumulam erro de deriva proporcional à distância percorrida. O Filtro de Madgwick
\cite{Madgwick2010} emerge como alternativa ao Filtro de Kalman Estendido (EKF) por
apresentar complexidade computacional $\mathcal{O}(1)$ por amostra, tornando-o
adequado para execução em dispositivos móveis de recursos limitados.

A abordagem híbrida BLE + PDR tem sido explorada na literatura para combinar a
precisão posicional dos beacons com a continuidade do PDR \cite{Jimenez2009}.
O presente trabalho diferencia-se por: (i) integrar a fusão sensorial a uma pilha
completa de roteamento geoespacial acessível; (ii) adotar arquitetura
\textit{offline-first} com sincronização por delta; e (iii) disponibilizar a
ferramenta como artefato reproduzível com implantação conteinerizada, aspectos
ausentes nos trabalhos correlatos.

% ============================================================
\section{Arquitetura do Sistema}
% ============================================================

O sistema é composto por três camadas principais, conforme ilustrado na
Figura~\ref{fig:arquitetura}: o \textbf{aplicativo Android}, o
\textbf{\textit{backend} FastAPI} e o \textbf{banco de dados geoespacial}.
A comunicação entre camadas utiliza GeoJSON como formato de serialização,
preservando a semântica geoespacial dos dados ao longo de toda a pilha.

\begin{figure}[ht]
\centering
\begin{verbatim}
  [Aplicativo Android]          [Backend FastAPI]
  BLE Scanner + PDR   <-HTTP->  /spaces/sync
  NavigationViewModel           /route/calculate
  Jetpack Compose               /analytics/ping
         |                            |
         |                    [PostgreSQL 16]
         |                    PostGIS + pgRouting
         +---- offline-first: grafo local ----+
\end{verbatim}
\caption{Visão geral da arquitetura do sistema de navegação indoor.}
\label{fig:arquitetura}
\end{figure}

\subsection{Banco de Dados Geoespacial}

O banco de dados é implementado em PostgreSQL 16 com as extensões PostGIS (dados
geoespaciais) e pgRouting (algoritmos de grafos). As entidades centrais são:

\begin{itemize}
  \item \textbf{Space:} representa o ambiente físico (ex.: \textit{shopping}, hospital);
  \item \textbf{Unit:} sala, loja ou departamento dentro do espaço, com atributo
        \texttt{is\_blocked} para indicar interdição temporária;
  \item \textbf{Beacon:} ponto de referência BLE com coordenadas geográficas conhecidas;
  \item \textbf{EdgeLine:} aresta do grafo de navegação com atributos de acessibilidade
        (\texttt{is\_accessible}, \texttt{restriction\_level}) que codificam escadas,
        rampas e elevadores;
  \item \textbf{POI:} ponto de interesse vinculado a uma unidade.
\end{itemize}

O esquema adota versionamento incremental via campo \texttt{version} na entidade
\texttt{Space}. O aplicativo consulta a versão antes de cada sessão e realiza
sincronização completa apenas quando detecta divergência, reduzindo o consumo de
dados em atualizações incrementais.

\subsection{\textit{Backend} — API REST}

O \textit{backend} é desenvolvido com FastAPI (Python 3.13+) e SQLAlchemy assíncrono
com \textit{driver} asyncpg. A Tabela~\ref{tab:endpoints} apresenta os principais
\textit{endpoints}.

\begin{table}[ht]
\caption{Endpoints principais da API de navegação indoor}
\label{tab:endpoints}
\centering
\begin{tabular}{|l|l|l|}
\hline
\textbf{Endpoint} & \textbf{Método} & \textbf{Função} \\
\hline
\texttt{/spaces/\{id\}/version} & GET  & Verificação de versão do mapa \\
\texttt{/spaces/\{id\}/sync}    & GET  & Sincronização GeoJSON completa \\
\texttt{/search}                & GET  & Busca textual de unidades e POIs \\
\texttt{/route/calculate}       & POST & Cálculo de rota com perfil de acessibilidade \\
\texttt{/analytics/ping}        & POST & Registro anônimo de presença \\
\texttt{/admin/spaces/publish}  & POST & Publicação de nova versão do mapa \\
\hline
\end{tabular}
\end{table}

O \textit{endpoint} \texttt{/route/calculate} recebe a posição de origem, o destino
e o perfil de acessibilidade do usuário, delegando o cálculo ao pgRouting via
algoritmo de Dijkstra sobre o grafo de \texttt{EdgeLines} filtrado pelo atributo
\texttt{restriction\_level}.

\subsection{Aplicativo Android}

O aplicativo é desenvolvido em Kotlin com Jetpack Compose como \textit{framework}
de interface, seguindo a arquitetura MVVM com \texttt{NavigationViewModel} como
ponto central de estado reativo (\texttt{StateFlow}). A comunicação com o
\textit{backend} utiliza Retrofit/OkHttp com serialização GeoJSON via
Kotlinx Serialization.

% ============================================================
\section{Posicionamento Indoor Híbrido}
% ============================================================

\subsection{Ancoragem por Beacons BLE}

Beacons BLE comerciais são distribuídos pelo espaço físico em posições geodésicas
previamente cadastradas no banco de dados. O aplicativo realiza varredura contínua
dos sinais BLE e identifica o beacon de maior RSSI (\textit{Received Signal Strength
Indicator}) como âncora de posição corrente. Ao detectar um beacon com RSSI superior
ao limiar configurável (padrão: $-70$ dBm), o sistema substitui a posição estimada
pelo PDR pelas coordenadas conhecidas do beacon, corrigindo o erro de deriva acumulado.

\subsection{Dead Reckoning Pedestre com Filtro de Madgwick}

O PDR estima o deslocamento do usuário a partir dos sensores inerciais do
\textit{smartphone}: acelerômetro, giroscópio e magnetômetro (configuração MARG —
\textit{Magnetic, Angular Rate and Gravity}). O Filtro de Madgwick \cite{Madgwick2010}
realiza a fusão desses sensores com complexidade computacional $\mathcal{O}(1)$ por
amostra, calculando o quaternion de orientação $\mathbf{q}$ de forma iterativa:

\begin{equation}
  \dot{\mathbf{q}} = \frac{1}{2}\mathbf{q} \otimes \boldsymbol{\omega} -
  \beta \frac{\nabla f}{\|\nabla f\|}
  \label{eq:madgwick}
\end{equation}

\noindent onde $\boldsymbol{\omega}$ é a velocidade angular medida pelo giroscópio,
$\beta$ é o ganho de correção (ajustável empiricamente) e $\nabla f$ é o gradiente
da função objetivo que minimiza o erro entre a aceleração medida e o vetor gravidade
estimado a partir da orientação atual.

O comprimento do passo é estimado dinamicamente a partir da variação da aceleração
vertical ao longo de uma janela temporal deslizante:

\begin{equation}
  L_{\text{passo}} = L_{\text{base}} \cdot \left(1 + k \cdot
  \frac{\sigma_a - \sigma_{\text{ref}}}{\sigma_{\text{ref}}}\right)
  \label{eq:step}
\end{equation}

\noindent onde $L_{\text{base}}$ é o comprimento de passo base calibrado por
usuário, $\sigma_a$ é o desvio padrão da aceleração na janela corrente e $k$
é um fator de escala empírico.

\subsection{Fusão Híbrida e Correção de Deriva}

Entre correções por beacon, a posição é propagada exclusivamente pelo PDR. A fusão
híbrida atua como um mecanismo de \textit{reset} periódico: cada beacon detectado
reinicia o acúmulo de erro inercial a partir de uma âncora de posição conhecida.
Essa estratégia dispensa redes densas de beacons — a cobertura não precisa ser
contínua, apenas suficiente para limitar o intervalo máximo de deriva.

% ============================================================
\section{Avaliação Experimental}
% ============================================================

\subsection{Configuração da Avaliação}

A validação preliminar foi conduzida nos corredores do Instituto de Informática da UFG
(Goiânia -- GO), ambiente indoor representativo com trechos de corredor, cruzamentos
e salas distribuídas em planta baixa irregular. Foram utilizados 4 beacons BLE
com distância máxima de 5 metros entre si, cobrindo o trajeto de avaliação.
O dispositivo móvel utilizado executava Android 7.0 ou superior, requisito mínimo
do aplicativo para acesso à API de sensores IMU via \texttt{SensorManager}.
O cálculo de rota é executado localmente no dispositivo via \texttt{RouteCalculator},
eliminando dependência de latência de rede durante a navegação.

Dado o caráter de prova de conceito desta avaliação, não foram realizadas medições
formais com gabarito físico ou número fixo de repetições controladas. Os resultados
apresentados a seguir são observacionais, obtidos durante sessões de uso do
aplicativo no ambiente descrito.

\subsection{Resultados Observacionais}

Durante as sessões de validação, o sistema demonstrou comportamento estável de
rastreamento ao longo dos corredores do Instituto de Informática. Os principais
resultados observados foram:

\begin{itemize}
  \item \textbf{Erro de posicionamento estimado:} a distância entre a posição
        exibida no aplicativo e a posição real do usuário foi estimada
        observacionalmente em aproximadamente 1,5 metros em condições normais
        de corredor, valor compatível com a literatura para fusão BLE + PDR
        em ambientes similares \cite{Faragher2015};
  \item \textbf{Estabilidade do rastreamento:} o Filtro de Madgwick manteve
        estimativa de orientação estável mesmo durante mudanças bruscas de
        direção e variações de ritmo de caminhada;
  \item \textbf{Correção por beacon:} a cada passagem por um beacon, o erro
        acumulado de deriva PDR foi visivelmente corrigido, confirmando o
        funcionamento do mecanismo de ancoragem;
  \item \textbf{Operação \textit{offline}:} após a sincronização inicial do
        mapa, a navegação funcionou completamente sem requisições adicionais
        ao \textit{backend}, validando a arquitetura \textit{offline-first}.
\end{itemize}

\subsection{Discussão e Limitações}

Os resultados indicam que a fusão BLE + PDR via Filtro de Madgwick é viável para
navegação indoor em ambientes de corredor com densidade moderada de beacons.
Identificam-se, contudo, as seguintes limitações desta avaliação: (i) a estimativa
de erro de posicionamento é observacional e não substitui medição formal por gabarito;
(ii) o erro de PDR cresce proporcionalmente em percursos longos sem cobertura de
beacon; (iii) a precisão do RSSI é sensível a obstruções físicas e variações de
tráfego humano; e (iv) a avaliação foi conduzida em ambiente de corredor único,
não cobrindo topologias com múltiplos andares. Uma avaliação quantitativa
controlada constitui trabalho futuro prioritário.

% ============================================================
\section{Disponibilidade e Demonstração}
% ============================================================

\subsection{Código-Fonte e Reprodutibilidade}

O sistema é disponibilizado integralmente sob a licença
\textbf{Creative Commons Attribution-ShareAlike 4.0 (CC BY-SA 4.0)},
incluindo \textit{backend} FastAPI, aplicativo Android e esquemas do banco de dados:

\begin{itemize}
  \item \textbf{Repositório:}
        \url{https://github.com/Jpborgesll/navegacao-indoor-srbc-2026}
  \item \textbf{Documentação e instalação:}
        \url{https://github.com/Jpborgesll/navegacao-indoor-srbc-2026#readme}
\end{itemize}

A reprodução do ambiente de execução é viabilizada por um arquivo
\texttt{docker-compose.yml} que provisiona automaticamente o banco de dados
PostgreSQL/PostGIS/pgRouting e o \textit{backend} FastAPI com um único comando
(\texttt{docker compose up --build}), eliminando dependências de configuração manual.

\subsection{Planejamento da Demonstração no SBRC}

A demonstração ao vivo será conduzida com os seguintes equipamentos:

\begin{itemize}
  \item \textbf{Smartphone Android} (API 26+) com sensores IMU e aplicativo instalado;
  \item \textbf{3 beacons BLE comerciais} (compatíveis com iBeacon/Eddystone)
        posicionados no espaço de demonstração;
  \item \textbf{Notebook} com Docker para execução local do \textit{backend} e
        do banco de dados;
  \item \textbf{Roteador Wi-Fi portátil} para comunicação entre dispositivos.
\end{itemize}

A demonstração consistirá em: (i) varredura e detecção dos beacons pelo aplicativo;
(ii) cálculo e exibição de rota acessível em tempo real; e (iii) rastreamento
contínuo de posição via fusão BLE + PDR durante o percurso do demonstrador.

% ============================================================
\section{Conclusão}
% ============================================================

Este trabalho apresentou um sistema de navegação indoor de baixo custo que combina
ancoragem por beacons BLE com Dead Reckoning Pedestre via Filtro de Madgwick em
uma arquitetura \textit{offline-first} completa. A validação preliminar nos corredores
do Instituto de Informática da UFG, com 4 beacons posicionados a até 5 metros entre
si, demonstrou rastreamento estável com erro estimado de aproximadamente 1,5 metros,
evidenciando a viabilidade técnica da abordagem em cenário indoor real com
infraestrutura de baixo custo.

Do ponto de vista de engenharia de software, a ferramenta contribui com uma
implementação reproduzível e de código aberto que integra fusão sensorial, roteamento
geoespacial acessível e analytics em conformidade com a LGPD, disponível para
uso e extensão pela comunidade de pesquisa.

Como trabalhos futuros, planeja-se: (i) validação em ambientes de maior escala e
complexidade topológica, incluindo múltiplos andares; (ii) calibração automática
do ganho $\beta$ do Filtro de Madgwick e do comprimento de passo via aprendizado
adaptativo; e (iii) extensão do módulo de analytics com detecção de zonas de
congestionamento em tempo real por clustering de trajetórias anonimizadas.

% ============================================================
%  REFERÊNCIAS
% ============================================================
\bibliographystyle{sbc}
\begin{thebibliography}{9}

\bibitem{Liu2007}
Liu, H., Darabi, H., Banerjee, P., and Liu, J. (2007).
Survey of wireless indoor positioning techniques and systems.
\textit{IEEE Transactions on Systems, Man, and Cybernetics, Part C},
37(6):1067--1080.

\bibitem{Zafari2019}
Zafari, F., Gkelias, A., and Leung, K. K. (2019).
A survey of indoor localization systems and technologies.
\textit{IEEE Communications Surveys \& Tutorials}, 21(3):2568--2599.

\bibitem{Faragher2015}
Faragher, R. and Harle, R. (2015).
Location fingerprinting with Bluetooth Low Energy beacons.
\textit{IEEE Journal on Selected Areas in Communications}, 33(11):2418--2428.

\bibitem{Jimenez2009}
Jiménez, A. R., Seco, F., Prieto, J. C., and Guevara, J. (2009).
A comparison of pedestrian dead-reckoning algorithms using a low-cost MEMS IMU.
In \textit{IEEE WISP 2009}, pages 37--42.

\bibitem{Madgwick2010}
Madgwick, S. O. H., Harrison, A. J. L., and Vaidyanathan, R. (2010).
Estimation of IMU and MARG orientation using a gradient descent algorithm.
In \textit{IEEE International Conference on Rehabilitation Robotics (ICORR)},
pages 1--7.

\end{thebibliography}

\end{document}
